\section{Going Further}
%% ═════════════════════════════════════════════════════════════════════════════

\begin{frame}{Going Further — pgAdmin}
  \begin{columns}[T]
    \column{0.55\textwidth}
      \begin{block}{\faDesktop\; pgAdmin}
        \textbf{pgAdmin} is the official open-source graphical desktop tool
        for managing PostgreSQL databases.

        \vspace{0.4em}
        \begin{itemize}
          \item Browse databases and tables visually
          \item Write and execute SQL queries in a GUI editor
          \item Inspect table structures, indexes, and constraints
          \item Monitor server activity and performance
          \item Export and import data easily
        \end{itemize}

        \vspace{0.4em}
        Download: \href{https://www.pgadmin.org/download/}{\texttt{pgadmin.org/download}}
      \end{block}

    \column{0.40\textwidth}
      \begin{exampleblock}{When to use \texttt{psql} vs pgAdmin}
        \vspace{0.2em}
        \begin{tabular}{p{1.4cm}p{2cm}}
          \textbf{psql} & Scripting, quick queries, remote servers \\[6pt]
          \textbf{pgAdmin} & Visual exploration, schema design, beginners \\
        \end{tabular}
      \end{exampleblock}

      \vspace{0.5em}
      \begin{alertblock}{\faLightbulb\; Tip}
        Start with pgAdmin to understand the structure visually, then move to \texttt{psql} for efficiency.
      \end{alertblock}
  \end{columns}
\end{frame}